\newcommand\odp[1]{\par\smallskip{\color{pomd}\bfseries #1}\par}
\odp{1. Wielokrotności}
Ćw. A: 12, 18, 24. Ćw. B: a) np. 0, 7, 14, 21, 28; b) 0, 14, 38, 150. Ćw. C: zielone 36 i 60 cm, czerwone 42 i 56 cm; inne – wielokrotności 12 i 14. Ćw. D: NWW(4, 10) = 20.\par
1. np. a) 9, 18, 27 b) 25, 50, 75 c) 120, 240, 360 d) 3000, 6000, 9000.
2. 1, 2, 3, 6, 9, 18.
3. a) 0, 18, 24, 51, 57, 72, 96 b) 0, 24, 72, 96 c) 0, 34, 51, 85.
4. a) 320 b) 280.
5. a) A = 48, C = 72 b) D = 261, F = 279, G = 288.
6. a) 20 b) 9.
7. P, F, P, P, F.
8. 4 (342 = 38 · 9).
9. a) 12 b) 24 c) 20 d) 347 e) 91 f) 30 g) 30 h) 12.\par
Super zagadka: kreski spotykają się co NWW(25, 20) = 100 cm, więc podwórko ma co najmniej 18 m – tak.

\odp{2. Dzielniki}
Ćw. A: a) 1, 2, 3, 4, 6, 8, 12, 24 b) 1, 3, 5, 15. Ćw. B: 1, 2, 4, 5, 10, 20. Ćw. C: a) 1, 2, 4, 7, 14, 28 b) 1; 2. Ćw. D: b) 1, 2, 3, 4, 6, 9, 12, 18, 36. Ćw. E: 4.\par
1. a) 1, 7 b) 1, 2, 4, 8 c) 1, 2, 4, 8, 16 d) 1, 3, 9, 27 e) 1, 2, 4, 5, 8, 10, 20, 40 f) 1, 2, 4, 8, 16, 32, 64.
2. P, P, F, P, F.
3. 1, 15, 30, 75, 150.
4. Drzwi 16 (1 + 2 + 4 + 8 + 16 = 31).
5. a) 200, 400, 600, 800 b) 150, 300, 450, 600, 750, 900 c) 250, 500, 750 d) 300, 600, 900.
6. a) 22 b) 36.
7. a) 1 + 2 + 3 = 6 – tak b) tak (1 + 2 + 4 + 7 + 14 = 28) c) nie (1 + 2 + 3 + 4 + 6 = 16).
8. np. 150; najmniejsza 120.
9. a) 6 b) 1 c) 5 d) 1 e) 1 f) 3 g) 1 h) 13 i) 1 j) 20 k) 1 l) 5.
10. 12 zestawów (po 2 jabłka i 3 gruszki).
11. a) 13, 26 b) 8, 16, 24 c) 5, 15, 25 d) np. 9 i 18, 9 i 27, 18 i 27.\par
Super zagadka: 21 (1 + 3 + 7 + 21 = 32).

\odp{3. Cechy podzielności}
Ćw. A: a) nie b) tak. Ćw. B: a) 340, 72, 64, 250, 900, 3400 b) 340, 55, 35, 105, 250, 175, 900, 815, 3400 c) 340, 250, 900, 3400 d) 900, 3400. Ćw. C: a) tak b) nie. Ćw. D: 436 i 9612. Ćw. E: nikt (np. 12). Ćw. F: 12, 27, 33, 51, 213 (obie grupy takie same). Ćw. G: 36 – tak, 9, tak; 47 – nie, 11, nie; 58 – nie, 13, nie; 702 – tak, 9, tak; 819 – tak, 18, tak; 974 – nie, 20, nie; 4042 – nie, 10, nie.\par
1. a) 20, 64, 488, 730, 1250, 2006, 4700, 98\,300, 2\,000\,500 b) 20, 85, 315, 730, 1250, 4700, 6205, 98\,300, 2\,000\,500 c) 20, 730, 1250, 4700, 98\,300, 2\,000\,500 d) 4700, 98\,300, 2\,000\,500.
2. Dwie ostatnie cyfry: 00, 25, 50 lub 75.
3. 1324, 516, 892, 21\,436.
4. 724, 936.
5. b) nie; 2028 i 2032 c) 1704, 1948, 2000, 2024 d) 1800, 1900, 2000 – przestępny tylko 2000.
6. 29 lutego 2028 r.
7. a) 135, 30\,012, 8127, 41\,022 b) 135, 8127, 41\,022.
8. a) np. 123 b) nie.
9. 7380.
10. przez 3: A, B, C, D, E; przez 9: B, C.
11. 1\,023\,456\,789 i 9\,876\,543\,210 – tak (suma cyfr 45).
12. a) 9 b) 11.
13. a) KOT, LOT.
14. tak; nie.
15. np. a) 1020, 10\,020 b) 1002, 10\,002 c) 1002, 10\,002.
16. a) np. 105 – tak b) 75 i 4605.
17. a) 2460, 1356, 91\,284 b) 2460, 1356, 91\,284 c) 2460, 74\,020.
18. przez 5: 2, 4, 1, 4, 0; przez 4: 1, 0, 2, 3, 2.
19. a) 1, 2, 0.\par
Super zagadka: 5 (2558 = 12 · 213 + 2).

\odp{4. Liczby pierwsze i złożone}
Ćw. A: 4, 2, 4, 2. Sito: 25 liczb.\par
1. np. 2316 – 2, 4785 – 5, 1818 – 2, 2721 – 3, 7413 – 3.
2. a) 101 i 997 b) 168 c) 830.
3. a) 1, 2003 b) 0, 2, 4, 5, 6, 8.
4. P, P, F, F, P.
5. 1.
6. 103, 827, 1993, 2027.\par
Super zagadka: 3; 4.

\odp{5. Rozkład na czynniki}
1. 18 = 2 · 3 · 3; 45 = 3 · 3 · 5; 72 = 2 · 2 · 2 · 3 · 3; 100 = 2 · 2 · 5 · 5; 126 = 2 · 3 · 3 · 7; 210 = 2 · 3 · 5 · 7.
2. a) 90 b) 56 c) 165.
3. 6; 36; 15; 90.
4. 3; 3.
